A family of subsets (``blocks'') $A_1,\dots,A_m\subseteq [n]:=\{1,\dots,n\}$ is a \emph{pairwise balanced block design} (PBD) if:
\begin{itemize}
\item every unordered pair $\{x,y\}\subseteq [n]$ with $x\neq y$ is contained in exactly one block $A_i$.
\end{itemize}
The problem asks: for all sufficiently large $n$, does there exist a \emph{nontrivial} PBD $A_1,\dots,A_m$ such that for every integer $t$,
\[
\#\{i:\;|A_i|=t\}=O(n^{1/2})
\quad\text{(uniformly in $t$)}?
\]

\emph{Ambiguity/nontriviality.} The single-block design $\{A_1=[n]\}$ is a PBD (every pair lies in that unique block). This is presumably the ``trivial'' design. The minimal corrected interpretation consistent with standard usage is:
\begin{quote}
``nontrivial'' means $m\ge 2$ (equivalently, not the single block $[n]$).
\end{quote}
Under this interpretation, every block must have size at most $n-1$.
